% Formalización LaTeX del documento INTERPRETACION_FIGURAS_FLUJOS_MATERIA_ACTIVA.md
% Compilar: pdflatex interpretacion_flujos_sigma_figuras.tex (dos pasadas si se usa hyperref)
\documentclass[11pt,a4paper]{article}

\usepackage[utf8]{inputenc}
\usepackage[T1]{fontenc}
\usepackage[spanish]{babel}
\usepackage{amsmath,amssymb,mathtools}
\usepackage{booktabs}
\usepackage{geometry}
\usepackage[hidelinks]{hyperref}
\usepackage{url}
\urlstyle{same}

\geometry{margin=2.5cm}

\newcommand{\R}{\mathbb{R}}
\newcommand{\vect}[1]{\mathbf{#1}}

\title{Flujos difusivos, gradientes y densidad $\sigma^{+}$:\\
marco termodinámico e interpretación de figuras\\
\large Modelo reacción--difusión tumor--sano--inmunidad (Allee)}
\author{}
\date{\today}

\begin{document}
\maketitle

\begin{abstract}
Se formaliza el continuo reacción--difusión en dos dimensiones para los campos
$(c,s,i)$, su encaje en el formalismo fenomenológico de la termodinámica del no
equilibrio (flujos y fuerzas) y la definición de la disipación difusiva
$\sigma_{\mathrm{diss,tot}}$ que el código grafica como $\sigma^{+}$.
Se describen los productos de salida del script
\texttt{visualize\_fluxes\_and\_entropy\_density.py} y una lectura en clave de
materia activa aplicada a cáncer, tejido sano e inmunidad.
\end{abstract}

\tableofcontents
\newpage

%---------------------------------------------------------------------
\section{Archivos generados y variables del modelo}
%---------------------------------------------------------------------

Para cada escenario, el script \path{nonequilibrium_termodynamics/visualize_fluxes_and_entropy_density.py}
escribe figuras en la subcarpeta \path{nonequilibrium_plots/}. Los nombres
típicos son:
\begin{center}
\begin{tabular}{@{}ll@{}}
\toprule
Archivo & Contenido \\
\midrule
\texttt{grad\_mag\_t\_$\langle t\rangle$.png} &
$\lVert\nabla c\rVert$, $\lVert\nabla s\rVert$, $\lVert\nabla i\rVert$ \\
\texttt{J\_mag\_t\_$\langle t\rangle$.png} &
$\lVert \vect{J}_c\rVert$, $\lVert \vect{J}_s\rVert$, $\lVert \vect{J}_i\rVert$ \\
\texttt{sigma\_plus\_t\_$\langle t\rangle$.png} &
$\sigma^{+}=\sigma_{\mathrm{diss,tot}}$ \\
\texttt{quiver\_Jc\_t\_$\langle t\rangle$.png} &
campo $\vect{J}_c$ sobre $c$ (opcional) \\
\bottomrule
\end{tabular}
\end{center}

Los campos $\phi_a\in\{c,s,i\}$ representan:
\begin{itemize}
  \item $c(\vect{x},t)$: densidad relativa de células tumorales;
  \item $s(\vect{x},t)$: densidad relativa de tejido sano;
  \item $i(\vect{x},t)$: actividad o densidad relativa del sistema inmune.
\end{itemize}

La opción \texttt{--sigma-csv} exporta la integral espacial de $\sigma^{+}$ en
función del tiempo.

%---------------------------------------------------------------------
\section{Contexto físico-matemático y termodinámico}
%---------------------------------------------------------------------

\subsection{Ecuaciones de campo (reacción--difusión)}

Sea $\Omega\subset\R^2$ el dominio espacial. Para $a\in\{c,s,i\}$,
\begin{equation}
\label{eq:rd}
\partial_t \phi_a = D_a \,\nabla^2 \phi_a + R_a(c,s,i),
\qquad D_a>0,
\end{equation}
donde $\nabla^2=\nabla\cdot\nabla$ es el laplaciano y $R_a$ agrupa las
contribuciones reactivas (Allee, competencia, inmunidad, etc.).

La forma \emph{conservativa} equivalente es
\begin{equation}
\label{eq:balance}
\partial_t \phi_a = -\nabla\cdot \vect{J}_a + R_a,
\qquad
\vect{J}_a = -D_a \nabla \phi_a .
\end{equation}
Aquí $\vect{J}_a$ es el flujo difusivo (ley de Fick isótropa lineal) y $R_a$ es
una fuente local no conservativa en general.

En la implementación numérica, $\phi_a$ se muestrea en una malla rectangular; los
gradientes en las figuras se aproximan por diferencias finitas con paso
$\Delta x$ deducido de \texttt{space\_size} y la resolución, en coherencia con
\path{termodynamics/calculate_thermodynamic_properties.py}.

\subsection{Termodinámica del no equilibrio: flujos y fuerzas}

En el formalismo de de Groot--Mazur (y relaciones de Onsager en el régimen
lineal), la densidad de producción de entropía local se escribe, en lo general,
como suma de productos bilineales flujo--fuerza:
\begin{equation}
\label{eq:sigmaS-general}
\sigma_S = \sum_\alpha J_\alpha X_\alpha ,
\end{equation}
bajo las hipótesis de validez del marco, con $\sigma_S\ge 0$ en situaciones
estándar.

Para difusión de especies a temperatura uniforme $T$, un término clásico es
\begin{equation}
\label{eq:sigmaS-diff}
\sigma_S = -\frac{1}{T}\sum_a \vect{J}_a\cdot\nabla\mu_a + \cdots
\end{equation}
donde $\mu_a$ es el potencial químico de la especie $a$. Si la constitutiva toma
la forma $\vect{J}_a=-D_a\nabla\mu_a$, entonces
\begin{equation}
\label{eq:Ja-mu}
-\vect{J}_a\cdot\nabla\mu_a = D_a \lVert\nabla\mu_a\rVert^2 \ge 0 .
\end{equation}

\subsection{Convención del módulo de visualización: $\mu_a\equiv\phi_a$}

El script de visualización no importa los $\mu_a$ obtenidos del funcional de
energía libre $F[c,s,i]$ del postproceso termodinámico. Trabaja con los propios
campos como variables termodinámicas efectivas del flujo:
\begin{equation}
\label{eq:Ja-phi}
\vect{J}_a = -D_a \nabla \phi_a .
\end{equation}

Se define la densidad de \emph{disipación difusiva} por especie y la total:
\begin{align}
\label{eq:sigma-diss-a}
\sigma_{\mathrm{diss},a} &:= D_a \lVert\nabla\phi_a\rVert^2 ,\\
\label{eq:sigma-plus}
\sigma^{+} = \sigma_{\mathrm{diss,tot}}
&:= \sum_a \sigma_{\mathrm{diss},a}
= D_c\lVert\nabla c\rVert^2 + D_s\lVert\nabla s\rVert^2 + D_i\lVert\nabla i\rVert^2 .
\end{align}
Las figuras \texttt{sigma\_plus\_t\_*.png} son mapas de $\sigma_{\mathrm{diss,tot}}(\vect{x},t)$,
coherentes con las series integrales guardadas por
\path{calculate_thermodynamic_properties.py}.

Introduciendo una temperatura efectiva constante $T_{\mathrm{eff}}$ (no
resuelta por el PDE), la analogía dimensional
\begin{equation}
\label{eq:sigmaS-supset}
\sigma_S \;\supset\; \frac{1}{T_{\mathrm{eff}}}\sum_a D_a\lVert\nabla\phi_a\rVert^2
\end{equation}
solo rescala amplitudes: no altera la forma espacial ni la evolución temporal
\emph{relativa} de $\sigma^{+}$.

\subsection{Distinción entre $\sigma_\mu$ y $\sigma_{\mathrm{diss,tot}}$}

Con potenciales $\mu_a$ derivados de un funcional $F[c,s,i]$,
\begin{equation}
\label{eq:sigma-mu}
\sigma_\mu := D_c\lVert\nabla\mu_c\rVert^2 + D_s\lVert\nabla\mu_s\rVert^2 + D_i\lVert\nabla\mu_i\rVert^2 .
\end{equation}
En general $\sigma_\mu \neq \sigma_{\mathrm{diss,tot}}$, salvo relaciones
particulares entre $(\mu_c,\mu_s,\mu_i)$ y $(c,s,i)$. Para interpretar las
figuras de este documento, el objeto relevante es $\sigma^{+}=\sigma_{\mathrm{diss,tot}}$
bajo \eqref{eq:Ja-phi}.

\subsection{Isotermia efectiva}

No existe campo $T(\vect{x},t)$ en el modelo implementado: el sistema es
\emph{isotermo efectivo}. Los campos $(c,s,i)$ son densidades poblacionales, no
energía interna. La comparación con la TdC es \emph{estructural} (balance
\eqref{eq:balance}, flujos, disipación por gradientes), no una medición literal
de entropía del tejido.

\subsection{Límites del proxy $\sigma^{+}$}

La densidad $\sigma^{+}$ no incluye:
\begin{itemize}
  \item una descomposición explícita $\sigma_{\mathrm{reacción}}(\vect{x},t)$ de
  las contribuciones entrópicas asociadas a $R_a$;
  \item intercambios con el exterior en un sistema biológico abierto.
\end{itemize}
Por tanto, $\sigma^{+}$ y su integral sobre $\Omega$ son observables
fenomenológicos del PDE: cuantifican el costo estructural de mezcla difusiva en
el continuo simulado.

\subsection{Energía libre $F$, potenciales y no reciprocidad}

El postproceso introduce $F[c,s,i]$ y $\mu_a\sim \delta F/\delta\phi_a$. El
jacobiano de las tasas,
\begin{equation}
\label{eq:jacobian-R}
A_{ab} := \frac{\partial R_a}{\partial \phi_b},
\end{equation}
admite la descomposición $A=S+N$ con $S$ simétrica y $N$ antisimétrica. Si
$N\neq 0$, el subsistema reaccional homogéneo no es integrable en general como
gradiente de un único potencial escalar en $(c,s,i)$; la dinámica global no está
garantizada como relajación pura de $F$. Esto no modifica las definiciones
\eqref{eq:Ja-phi}--\eqref{eq:sigma-plus}, pero acota el lenguaje al vincular
las figuras con ``descenso de $F$'' o equilibrio químico local.

\subsection{Correspondencia figura--formalismo}

\begin{center}
\begin{tabular}{@{}p{3.2cm}p{4.5cm}p{5.5cm}@{}}
\toprule
Figura & Objeto & Rol (analogía TdC) \\
\midrule
\texttt{grad\_mag} &
$\lVert\nabla\phi_a\rVert$ &
pendiente espacial; fija, con $D_a$, flujo y $\sigma_{\mathrm{diss},a}$ \\
\texttt{J\_mag} &
$\lVert\vect{J}_a\rVert = D_a\lVert\nabla\phi_a\rVert$ &
flujo difusivo \\
\texttt{sigma\_plus} &
$\sigma_{\mathrm{diss,tot}}$ &
proxy de parte difusiva de disipación / $\sigma_S$ \\
\texttt{quiver\_Jc} &
$\vect{J}_c$ sobre $c$ &
dirección del transporte de $c$ \\
\bottomrule
\end{tabular}
\end{center}

%---------------------------------------------------------------------
\section{Interpretación de cada tipo de figura}
%---------------------------------------------------------------------

\subsection{Gradientes: \texttt{grad\_mag\_t\_*.png}}

Se grafica $\lVert\nabla\phi_a\rVert$ para $a\in\{c,s,i\}$. Valores altos
indican frentes e interfaces; valores bajos, regiones casi uniformes. Los tres
paneles comparten escala máxima (\texttt{share\_vmax}), lo que permite comparar
heterogeneidad relativa entre compartimentos antes de ponderar por $D_a$.

\subsection{Magnitudes de flujo: \texttt{J\_mag\_t\_*.png}}

Se muestra $\lVert\vect{J}_a\rVert = D_a\lVert\nabla\phi_a\rVert$. Un gradiente
grande con $D_a$ pequeño puede producir flujo débil, y viceversa. La comparación
con \texttt{grad\_mag} distingue ``mucho contraste'' de ``mucho transporte''.

\subsection{Densidad $\sigma^{+}$: \texttt{sigma\_plus\_t\_*.png}}

Mapa de $\sigma_{\mathrm{diss,tot}}$ \eqref{eq:sigma-plus}. Los máximos locales
suelen coincidir con frentes donde la mezcla difusiva es intensa. La integral
discreta $\sum_{ij}\sigma^{+}(\vect{x}_{ij})\,(\Delta x)^2$, reportada por el
script, resume la disipación difusiva global en el instante $t$.

\subsection{Quiver de $\vect{J}_c$: \texttt{quiver\_Jc\_t\_*.png}}

Flechas (submuestreadas) con la dirección de $\vect{J}_c=-D_c\nabla c$ sobre el
fondo $c(\vect{x},t)$; complementa \texttt{J\_mag} con orientación en el plano.

\subsection{Lectura temporal}

Comparar distintos $t$ con el mismo tipo de figura revela migración de frentes,
fragmentación u homogeneización, y la evolución de la integral de $\sigma^{+}$.
Si $\lVert\nabla c\rVert$ es grande pero $\lVert\vect{J}_c\rVert$ no, conviene
revisar $D_c$ en los parámetros del escenario.

%---------------------------------------------------------------------
\section{Materia activa: analogías para tumor, tejido sano e inmunidad}
%---------------------------------------------------------------------

La materia activa designa, en física, sistemas de muchos agentes que consumen
energía y generan fases y patrones fuera del equilibrio térmico pasivo. El
presente modelo es un \emph{continuo} reacción--difusión, no un enjambre
microscópico de partículas autopropulsadas; no obstante, varias nociones
iluminan la lectura de las figuras.

\paragraph{Tres especies acopladas.}
Las fuentes $R_a$ mantienen el sistema alejado del equilibrio trivial: $c$
(invasión, Allee), $s$ (competencia, límites con $c$ e $i$), $i$ (retroalimentación
inmune, p.\,ej.\ controles tipo Hill en $R_i$). Las figuras muestran campos
efectivos; la ``actividad'' microscópica está promediada en $R_a$.

\paragraph{Frentes e interfaces.}
Los frentes concentran intercambio y disipación efectiva: \texttt{grad\_mag}
marca bordes; \texttt{sigma\_plus} suele destacar esas regiones como hotspots de
$\sigma_{\mathrm{diss,tot}}$.

\paragraph{Corrientes.}
Solo aparece explícitamente la corriente difusiva $\vect{J}_a$. Las reacciones
moldean los campos en el tiempo pero no se representan como campos vectoriales
adicionales en el plano: conviene separar \emph{transporte difusivo} y
\emph{dinámica reaccional local} en la redacción científica.

\paragraph{Lenguaje entrópico.}
Para tesis o artículos, conviene declarar isotermia efectiva, usar
$\sigma_{\mathrm{diss,tot}}$ como proxy de la parte difusiva de la producción
entrópica, y mencionar la no reciprocidad en $A$ si se argumenta con $\mu_a$
derivado de $F$.

%---------------------------------------------------------------------
\section{Referencias en el repositorio}
%---------------------------------------------------------------------

\begin{itemize}
  \item \path{nonequilibrium_termodynamics/MARCO_FISICO_MATEMATICO.md}
  \item \path{nonequilibrium_termodynamics/visualize_fluxes_and_entropy_density.py}
  \item \path{termodynamics/calculate_thermodynamic_properties.py}
  \item \path{nonequilibrium_termodynamics/reciprocity_jacobian_analysis.py}
  \item \path{Allee/PIPELINE_EJECUCION_Y_FISICA.md}
\end{itemize}

\vspace{1em}
\noindent\textit{Las ecuaciones \eqref{eq:Ja-phi}--\eqref{eq:sigma-plus} coinciden con la implementación del script de visualización y con el MARCO \S4.1.}

\end{document}
